\section{The RI/DF path: smaller derivative memory}
\label{sec:ri-df}

The nuclear-derivative path of the dense four-center integrals allocates
a \texttt{nctr}~$\times$~\texttt{nao}$^4$ array; at \texttt{nao}~$=96$
that is 55~GB, and at \texttt{nao}~$=175$ it is 2~TB.  The RI/DF layer
replaces those integrals in the correlation stage by two three-center
factors and a metric, and exposes a derivative path whose working set is
\texttt{nao}$^2$~$\times$~\texttt{naux} per atom.  This chapter runs
that layer on H$_2$O/6-31G with CAS(2,2) and ncore~$=4$, checks the
assembly and the CI energy against the dense route, and verifies one
derivative against a central difference.  The manual's RI chapter
carries the formulas, the full API and the error budget.

\subsection{Build the factors}

\begin{lstlisting}[language=Python]
import jax
jax.config.update('jax_enable_x64', True)

from pyscf import gto, scf
from diffmrc.core import ri

mol = gto.M(atom='O 0 0 0; H 0.7557 0 0.5880; H -0.7557 0 0.5880',
            basis='6-31g', unit='Angstrom', verbose=0)
mf = scf.RHF(mol).run(verbose=0)
auxmol, label = ri.make_auxmol(mol)
factors = ri.ri_factors(mol, auxmol)

print(label, factors.nao, factors.naux, factors.cond)
# def2-universal-jkfit 13 113 2.47e+07
\end{lstlisting}

\texttt{make\_auxmol} tries the fitting bases in
\texttt{ri.AUX\_BASIS\_PRIORITY} and reports which one it used; the
returned label is what should be quoted with any RI number.  On this
system the metric condition number is $2.5\times10^{7}$ and no direction
is dropped at the default \texttt{rcond}~$=10^{-12}$.

\subsection{Compare the assembled tensor with the dense one}

\begin{lstlisting}[language=Python]
import numpy as np

eri_ri = ri.ri_eri_ao(mol, auxmol, factors=factors)
eri_dense = mol.intor('int2e', aosym='s1')

bf = factors.b.reshape(-1, factors.naux)
ref = bf @ np.linalg.solve(factors.s, bf.T)
print('identity', np.abs(eri_ri - ref.reshape(eri_ri.shape)).max())
print('rms     ', np.sqrt(np.mean((eri_ri - eri_dense) ** 2)))
print('max     ', np.abs(eri_ri - eri_dense).max())
\end{lstlisting}

The identity check compares the two-step GEMM assembly with an explicit
metric solve; it is the implementation error and should print at the
$10^{-12}$ level.  The
direct comparison with the dense tensor is the RI approximation itself
and is four to eight orders of magnitude larger.  On this run:

\begin{lstlisting}
identity 1.068e-12
rms      8.317e-06
max      7.967e-04
\end{lstlisting}

\subsection{RI-CASCI}

The MO path transforms the three-center tensor once and forms the
four-center MO integrals by one GEMM over the auxiliary index:

\begin{lstlisting}[language=Python]
from diffmrc.core.casci import cs_casci

c = mf.mo_coeff
h1_mo, eri_mo_dense = mo_integrals(c, mf.get_hcore(), eri_dense)
b_mo, eri_mo_ri = ri.ri_eri_mo(c, factors.b, factors.jhalf)

kw = dict(ncore=4, nact=2, nele_act=2, return_ci=True)
e_dense, ci_dense = cs_casci(h1_mo[:6, :6], eri_mo_dense[:6, :6, :6, :6], **kw)
e_ri, ci_ri = cs_casci(h1_mo[:6, :6], eri_mo_ri[:6, :6, :6, :6], **kw)
print(float(e_ri) - float(e_dense))
# +1.054e-05
\end{lstlisting}

\texttt{mo\_integrals} is the four-step two-operand transform defined in
the example script.  The RI shift of the CASCI energy is
$1.1\times10^{-5}$~Eh (0.007~kcal/mol), and the CI vectors overlap to
ten digits.  At this size the shift is dominated by the fitting-basis
error, not by the contraction.

\subsection{Check a derivative against central differences}

\begin{lstlisting}[language=Python]
def displaced(mol, atom, axis, step):
    coords = mol.atom_coords(unit='Bohr').copy()
    coords[atom, axis] += step
    symbols = [mol.atom_pure_symbol(i) for i in range(mol.natm)]
    return gto.M(atom=[(s, c) for s, c in zip(symbols, coords)],
                 basis=mol.basis, unit='Bohr', verbose=0)

step = 1e-4
for atom in range(mol.natm):
    d3 = ri.ri_three_center_deriv(mol, auxmol, atom)
    for axis in range(3):
        mp = displaced(mol, atom, axis, step)
        mm = displaced(mol, atom, axis, -step)
        ap, _ = ri.make_auxmol(mp)
        am, _ = ri.make_auxmol(mm)
        fd = (ri.ri_three_center(mp, ap)
              - ri.ri_three_center(mm, am)) / (2 * step)
        print(atom, axis,
              np.abs(d3[axis] - fd).max() / np.abs(fd).max())
\end{lstlisting}

The auxiliary molecule is rebuilt at each displaced geometry so that the
fitting basis follows the nuclei.  The worst relative deviation over all
atoms and axes is

\begin{lstlisting}
worst relative deviation   3.641e-09  (atom 0, axis 0)
\end{lstlisting}

The same comparison for the assembled ERI,
\texttt{ri\_eri\_deriv} against a central difference of
\texttt{ri\_eri\_ao}, gives $3.7\times10^{-8}$ at the same step.  Part
of that residual is the finite-difference truncation amplified by the
metric inverse square root; the manual's error budget separates the two
contributions.

\subsection{The NEVPT2 comparison}

\texttt{examples/ri\_nevpt2.py} feeds the RI MO integrals to
\texttt{core.nevpt2.impl.nevpt2\_e2\_full2} and compares the correction
with the dense route at the same orbitals.  On H$_2$O/6-31G,
CAS(2,2), ncore~$=4$:

\begin{lstlisting}
e2 dense              -0.1276536999 Eh
e2 RI, same CI        -0.1276366317 Eh
e2 RI, end to end     -0.1276366150 Eh
RI - dense, same CI   +1.707e-05 Ha
RI - dense, end to end +1.708e-05 Ha
\end{lstlisting}

The three numbers separate the integral approximation (same CI) from
the small CI shift caused by it (end to end).  The absolute correction
is that of a fixed-orbital calculation, not a production NEVPT2 value;
what the example validates is the RI-minus-dense difference.

\subsection{Running the examples}

From the repository root, with the package importable:

\begin{lstlisting}[language=bash]
$ python examples/ri_df.py
$ python examples/ri_nevpt2.py
# if the package is not installed:
$ PYTHONPATH=src python examples/ri_df.py
\end{lstlisting}

Both scripts set \texttt{JAX\_PLATFORMS=cpu} if it is unset and require
only \texttt{pyscf}, \texttt{jax}, \texttt{numpy} and \texttt{diffmrc}.
On the reference machine, \texttt{ri\_df.py} takes 6.4~s of script time
(8.0~s wall, 0.9~GB peak RSS) and \texttt{ri\_nevpt2.py} takes 36~s of
script time (40~s wall, 1.6~GB peak RSS); the NEVPT2 part is dominated
by the first JAX compilation.
